
\chapter{Business Understanding}
\label{chpt:BU}
\graphicspath{{Chapter1/Chapter1Figs/PNG/}{Chapter1/Chapter1Figs/PDF/}{Chapter1/Chapter1Figs/}}



\section{Business Objectives}

In business terms, state what to you want to accomplish when mining this data.


\section{Assess the situation}

If there are any issues with access to data or time constraints, mention then here. Also details risks if any. Otherwise leave out this section.


\section{Data Mining Objectives}

State your technical objectives for mining the data.



\section{Project plan}
See suggested project plan on Moodle and adjust to suit your own project.

An example of a table is shown in \ref{IntFactors} 


\begin{table} [!t]\footnotesize    %!t here should force the table to appear at the top of the page
\centering
\captionsetup{font=small}
\captionsetup{aboveskip=-0.2\normalbaselineskip}
\caption[Defining broad factors of intelligence]{Defining broad factors of intelligence}
\label{IntFactors}

\begin{tabular}{p{3.5cm} l p{8cm} }\\   
\noalign{\smallskip}\hline\noalign{\smallskip}

Factor & Symbol & Description \\
\noalign{\smallskip}\hline\noalign{\smallskip}
Fluid Intelligence	&Gf 	&Ability to solve problems independently of knowledge learned.\\
\noalign{\smallskip}
Crystallised\newline intelligence&	Gc	&Acquiring and organising knowledge and skills, and ability to use such knowledge in solving problems.\\
\noalign{\smallskip}
Visual processing&	Gv&	Ability to process and analyse visual information\\
\noalign{\smallskip}
Auditory Processing	&Ga&	Ability to process and analyse auditory information\\
\noalign{\smallskip}
Processing Speed&	Gs&	Ability to perform automatic cognitive tasks quickly (measure in minutes)\\
\noalign{\smallskip}
Reaction Time / \newline Decision speed&	Gt&	Speed at which an individual can react to a stimulus, or make decisions (measure in seconds).\\
\noalign{\smallskip}
Short-Term Memory	&Gsm&	Ability to hold information with immediate awareness and use it again within a few seconds\\
\noalign{\smallskip}
Long-Term Retrieval&Glr&Ability to store and retrieve information over a longer period of time.\\
\noalign{\smallskip}
Quantitative\newline Knowledge&Gq	&Ability to understand quantitative concepts and relationships, and work with numeric symbols. This is a measure of mathematical knowledge acquired, as distinct from mathematical reasoning (Gf). \\
\noalign{\smallskip}
Reading-Writing	&Grw&	Basic reading and writing skills (considered by Cattell-Horn to be part of Gc)\\
\hline
\end{tabular}
\par\medskip
\end{table}	

Different citation styles: Cognitive ability tests were originally developed to identify low academic achievers \citep{Jensen1981, Munzert1980}. The first such test measured general cognitive intelligence, \emph{g}, as identified by \citet{Spearman1904, Spearman1927}. 



